\documentclass[11pt]{article}
\usepackage[a4paper,margin=25mm]{geometry}
\usepackage{amsmath,amssymb,amsthm,hyperref}
\newtheorem{theorem}{Theorem}
\newtheorem{lemma}{Lemma}
\newcommand{\PP}{\mathcal P}
\title{Two polynomial-size dice in a completely finite fair family}
\author{Complete proof; independently reviewed}
\date{30 September 2026}
\begin{document}\maketitle
\noindent\textbf{Scope.} This proof combines the reviewed real step-rank
construction and positive completion with a rational smoothing argument
and the companion clustered lifting lemma. All parts have passed full independent mathematical reviews by the root,
rational-design, and independent-direction agents. This is not formal
verification. The other dice may be extremely large.

\begin{theorem}\label{main}
There are absolute constants $C,D$ such that for every $n\ge3$ a
completely finite permutation-fair family of $n$ equiprobable dice
has two specified dice with at most $Cn^D$ faces each.
The other $n-2$ dice can have a common finite face count.
The proof below permits $D=162$; no attempt is made to optimize it.
\end{theorem}

Put $m=n-2$. The two reviewed inputs are
\path{../rational_designs/asymptotic/two_discrete_step_compatibility.tex}
and
\path{../independent_directions/two_exceptional_positive_completion.tex}.
They supply distinct nodes $0<b_1<\cdots<b_K<1$ and increasing rational
ranks $0<r_1<\cdots<r_K<1$, with
\begin{equation}\label{input}
 \frac1K\sum b_q^j=\frac1{j+1}\quad(0\le j\le m+1),\qquad
 \frac1K\sum r_qb_q^j=\frac1{j+2}\quad(0\le j\le m),
\end{equation}
where $K=O((m+1)^4)$ and the $r_q$ have a common denominator
$L=O((m+1)^{160})$. Their closeness guarantees a positive completion
$f_0$, with $1/2\le f_0\le3/2$, on the point-node gaps.

We also use a property proved in the step-rank construction: with $r$
fixed, the Jacobian in $b$ of the $m+1$ ordinary moment equations and
the $m$ nonconstant rank-weighted equations has row rank $2m+1$.
Indeed, before the small Newton correction its rows evaluate
$\PP_m+(r-b)\PP_{m-1}$, whose Gram matrix is positive by the Legendre
conditioning lemma. The fixed-right-inverse Newton map in that proof is a contraction
on its affine correction subspace. At its solution, therefore, the
right inverse composed with the new Jacobian is invertible on that
subspace. This directly preserves row rank $2m+1$; no further
quantitative lower bound is needed.

\begin{lemma}[Rational positive smoothing of the discrete variable]
There are disjoint rational closed intervals $J_q\subset(0,1)$ and a
positive affine density on each $J_q$, zero elsewhere, defining a
probability measure $\beta$, such that
\begin{align}
 \beta(J_q)&=1/K,\label{mass}\\
 \int x^j\,d\beta(x)&=1/(j+1) &&(0\le j\le m+1),\label{bmom}\\
 \sum_q r_q\int_{J_q}x^j\,d\beta(x)&=1/(j+2)
                                      &&(0\le j\le m).\label{weighted}
\end{align}
The conditional density in every box, rescaled to $[-1,1]$, is between
$1/4$ and $3/4$. All density coefficients are rational. The boxes may
be made arbitrarily close to the points $b_q$ and arbitrarily short.
\end{lemma}
\begin{proof}
Let $\varepsilon\downarrow0$ through positive rational numbers and
choose rational centers $c_q$ with $|c_q-b_q|\le\varepsilon^3$.
Set $J_q=[c_q-\varepsilon,c_q+\varepsilon]$; these are disjoint and
interior for small enough $\varepsilon$. On $J_q$ use density
\[
 \frac1{2K\varepsilon}
       \left(1+a_q\frac{x-c_q}{\varepsilon}\right).
\]
Its box mass is $1/K$ for all $a_q$.
At $a=0$, the residual in the $2m+1$ nonconstant constraints
\eqref{bmom}--\eqref{weighted} is $O_m(\varepsilon^2)$ by symmetric
Taylor expansion and \eqref{input}.
The coefficient matrix for the variables $a_q$, divided by
$\varepsilon$, converges to one third of the full-rank Jacobian
specified above. This follows from
\[
 \frac12\int_{-1}^1 yP(c_q+\varepsilon y)\,dy
     =\frac\varepsilon3P'(b_q)+O_m(\varepsilon^3).
\]
Choose a nonsingular square minor of that Jacobian, set the other
$a_q$ to zero, and solve the corresponding linear system exactly.
It gives $a_q=O_m(\varepsilon)$, hence $|a_q|\le1/2$ for sufficiently
small boxes. All entries of the actual system and its targets are
rational, so all $a_q$ are rational. The two degree-zero constraints
are already automatic from the box masses and $K^{-1}\sum r_q=1/2$.
This proves the lemma.
\end{proof}

\begin{lemma}[Positive rational completion in the complementary gaps]
For all sufficiently short boxes in the preceding lemma, there is a
rational piecewise polynomial density $f$, supported on their complement,
with mass $r_{q+1}-r_q$ on the $q$th complementary gap, where
$r_0=0,r_{K+1}=1$. On every such gap, $1/4\le f\le7/4$.
If $A$ has density $f$, $B$ has law $\beta$, and $m$ further variables
are independent uniforms, all complete orders are uniform.
\end{lemma}
\begin{proof}
Write $G_\varepsilon=[0,1]\setminus\bigcup_qJ_q$ and define
\[
 S_j^\varepsilon(a)=\int_{b>a}b^j\,d\beta(b),\qquad 0\le j\le m.
\]
On each component of $G_\varepsilon$ these are constant. Let
\[
 V_\varepsilon=\operatorname{span}\{a^j:0\le j\le m\}
 +\operatorname{span}\{a^rS_s^\varepsilon(a):r+s\le m\},
 \qquad S_\varepsilon=\operatorname{span}\{1,S_0^\varepsilon,\ldots,S_m^\varepsilon\},
\]
where the spaces are restricted to $G_\varepsilon$. Let $P_\varepsilon$
be averaging separately on its complementary gaps.

The matrix
\[
 M_{q,s}^\varepsilon=\int_{J_q}b^s\,d\beta(b),\qquad 0\le s\le m,
\]
has full column rank for sufficiently small boxes, since it converges
to the scaled Vandermonde matrix $K^{-1}(b_q^s)$.
Consequently
\begin{equation}\label{kernel}
 V_\varepsilon\cap\operatorname{ran}P_\varepsilon=S_\varepsilon.
\end{equation}
To see this directly, write a function in $V_\varepsilon$ as
$p(a)+\sum_{r+s\le m}c_{r,s}a^r S_s^\varepsilon(a)$.
If it is constant in each gap, the coefficient of $a^r$, for every
$r\ge1$, vanishes in all gaps. Subtracting consecutive gaps yields
$\sum_s c_{r,s}M_{q,s}^\varepsilon=0$ for all $q$.
The full-rank property, also for each initial set of columns, gives
$c_{r,s}=0$ and then the coefficient of $a^r$ in $p$ is zero.
Only the displayed step space remains. The same argument, including
$r=0$, proves independence of the indicated polynomial/step basis.

Let $g_\varepsilon$ be constant on each gap with the prescribed mass
$\gamma_q=r_{q+1}-r_q$. Seek a zero-gap-mean polynomial correction
$h_\varepsilon$ in
$W_\varepsilon=(I-P_\varepsilon)V_\varepsilon$ such that
\begin{equation}\label{riesz}
 \int_{G_\varepsilon}h_\varepsilon F
   =\int_0^1F(a)\,da-\int_{G_\varepsilon}g_\varepsilon F
                         \qquad(F\in V_\varepsilon).
\end{equation}
For the full-interval integral, use the functions before restriction.
The right side annihilates $S_\varepsilon$: the constant is immediate,
and Fubini, the gap masses, \eqref{bmom}, and \eqref{weighted} give
\[
 \int_{G_\varepsilon}g_\varepsilon S_j^\varepsilon
   =\sum_q r_q M_{q,j}^\varepsilon
   =\frac1{j+2}
   =\int_0^1 S_j^\varepsilon(a)\,da.
\]
Thus \eqref{kernel} makes \eqref{riesz} a well-defined Riesz problem
on $W_\varepsilon$.

Here is the needed continuity without any uniform lower bound on box
widths. A basis of $W_\varepsilon$ is given by applying
$I-P_\varepsilon$ to $a^r$ ($1\le r\le m$) and to
$a^r S_s^\varepsilon$ ($r\ge1$, $r+s\le m$).
The preceding full-rank argument proves its independence. Its Gram
matrix and the right-hand vector in \eqref{riesz} converge as
$\varepsilon\to0$ to those of the point-node problem in the reviewed
positive-completion proof. That limiting Gram matrix is positive
definite. Therefore the coefficients of $h_\varepsilon$ converge to
those of its point-node representative. After identifying each moving
gap affinely with its original gap, all functions involved are
polynomials of bounded degree with convergent coefficients. Convergence
is uniform on each closed gap. The limiting completed density is the
reviewed $f_0\in[1/2,3/2]$. Hence
$f=g_\varepsilon+h_\varepsilon\in[1/4,7/4]$ on all gaps for sufficiently
small boxes. All Gram entries and right-hand entries are rational:
box and gap endpoints, ranks, and polynomial density coefficients are
rational. Thus the coefficients of $f$ are rational as well.

By \eqref{riesz}, $f$ integrates every test in $V_\varepsilon$ exactly
as Lebesgue measure does. For $r+s\le m$,
\[
 \mathbb E[A^rB^s\mathbf1_{A<B}]
 =\int f(a)a^rS_s^\varepsilon(a)\,da
 =\frac1{r+1}\int b^{r+s+1}\,d\beta(b)
 =\frac1{(r+1)(r+s+2)}.
\]
The ordinary moments and the opposite triangle follow in the same way
as in the reviewed point-node completion. These are all the polynomial
chamber moments needed for $m$ other uniform variables, proving full
order fairness.
\end{proof}

\begin{proof}[Proof of Theorem~\ref{main}]
We have rational polynomial densities for $A$ and $B$ on alternating,
disjoint interval interiors. Their cluster masses have denominators
$L$ and $K$, respectively. The conditional densities, after rescaling
a cluster to $[0,1]$, are bounded above and below by absolute positive
constants: for $A$ use the bounds $1/4,7/4$ and its cell average;
for $B$ use its affine density bounds.

Choose an integer $T=\lceil C(m+1)^2\rceil$ with sufficiently large absolute $C$.
On each $B$ box take an equal-weight quadrature with exactly $T$ nodes.
On an $A$ gap of mass $a_q/L$, take one with exactly $Ta_q$ nodes.
The Gilboa--Peled every-integer-node-count theorem applied at degree
$2m+2$ allows these choices, since $a_q\ge1$. The standard interior
Vandermonde regularization preserves moments through degree $m$ and
makes all nodes distinct and strictly interior: there are at least
$m+2$ distinct support points in the higher-degree rule, hence at least
$m$ interior pivots, and the implicit-function argument removes
endpoints and collisions. This argument is detailed in the reviewed
one-exceptional proof and its every-size corollary.

Globally, the $B$ rule has $KT=O((m+1)^6)$ equally likely nodes and the
$A$ rule has $LT=O((m+1)^{162})$ equally likely nodes. Their moments
through degree $m$ on every individual cluster are rational, being the
integrals of rational polynomial densities over rational intervals.
Every cluster contains at least $m^2$ distinct nodes. The new hybrid
model is still fully fair: on a product of an $A$ cluster and a $B$
cluster, their relative order is fixed and every conditional old-order
kernel is polynomial of degree at most $m$ in each specified variable.
The cellwise quadratures preserve its integral.

Shrink the cluster intervals slightly around these finitely many
interior nodes to make their closures disjoint. The theorem in
\path{clustered_multi_exceptional_lifting.tex} now applies with $k=2$.
It retains the two displayed face counts and replaces the $m$ uniforms
by finite dice with one common face count. This proves the result.
\end{proof}
\end{document}
